% Fallbacks and common problems of the Day 11 lab. The rows come from
% Labs/day11/README.md. Nobody tested the fixes on the hardware.
\begin{frame}{If a part does not work, and common problems}
  \footnotesize
  \renewcommand{\arraystretch}{1.1}
  \setlength{\tabcolsep}{4pt}
  \begin{tabular}{@{}L{0.42\textwidth}L{0.53\textwidth}@{}}
    \toprule
    \textbf{Problem} & \textbf{Fix or fallback} \\
    \midrule
    \code{ERROR: cannot open the stream}
      & Check the log of MediaMTX and \code{ss -tln}. Use the address in
        place of \code{pi-NN.local}. \\
    An error for the camera
      & A second program uses the camera, for example a Day 8 script. Stop
        it. \\
    The camera does not work
      & Send a test image to the path \code{test} (HW-06, step 8) \\
    \code{ERROR: no file detector.py}
      & Give the folder: \code{--day08 <path>/day08/pi} \\
    \code{No module named 'ultralytics'}
      & Start the script with \code{\$PY}, the Day 8 environment \\
    Grey areas in the image for a short time
      & Lost packets with UDP. Use TCP, or write it in the report. \\
    The latency grows during a run
      & Use the reader of the lab (no \code{--in-order}) \\
    \code{mjpeg\_rate.py} waits, the browser shows the stream
      & The sketch serves one reader at a time: close the browser tab \\
    The XIAO does not connect, or \code{Camera init failed}
      & 2.4 GHz and the password. \code{OPI PSRAM}, the camera
        connector. \\
    A task is not complete in time
      & Use the file of \code{solutions/}, and write this in the report \\
    \bottomrule
  \end{tabular}
  \vskip 0.2em
  \small
  The README of the lab has the complete list.
\end{frame}
